\documentclass[11pt]{article}
\usepackage{../shared/luxcover}
\usepackage{amsmath,amssymb}
\usepackage{hyperref}
\usepackage{booktabs}
\usepackage{enumitem}
\usepackage{xcolor}

\title{Lux Silver}
\author{Zach Kelling\\\textit{Lux Industries, Inc.}\\\texttt{research@lux.network}}
\date{April 2024}

\begin{document}

\luxcoverpage

\begin{abstract}
Invest in the future of Silver
\end{abstract}

\tableofcontents
\newpage


\section{Abstract}

Lux Silver aims to provide a convenient, secure and easy to use investment alternative for modern investors interested in holding Silver.


Our mission is to become one of the fastest growing Silver companies by utilizing Blockchain technology and modern strategies.


The Silver Cliff project, managed by Viscount Mining, is one of the most promising silver mines in North America. Lux Silver brings this valuable asset into the digital age, offering a unique investment opportunity. The market for silver is robust and growing, driven by its use in various industries, including electronics, solar energy, and jewelry. The introduction of blockchain technology into the silver market through the Lux Silver project is expected to increase accessibility and liquidity, potentially driving further growth in the market.


\section{Highlights}

The World's Largest DeFi Physical Silver Fund -- The only physical Silver fund to provide a convenient way to stake or liquidate your tokenized physical silver assets.


Experienced Mining Team Lux Silver's partnership with Viscount Mining, a reputable mining company with a portfolio of silver and gold properties in the Western United States, ensures the reliability and credibility of the silver resources backing each token.


Low Risk -- Tokens are backed by physical Silver making it a very low risk investment opportunity, given enormous global demand.


Convenient Way to Own Physical Silver -- It's easy to buy, stake or sell silver on the Lux marketplace.


\section{Key Benefits}

\subsection{For retail and institutional investors}

Securely assessed, real verified value / NI 43-101 indicated amount of over 10 million Oz of silver, over 43 million ounces of future production available.


Easy to trade with others / buy sell or stake on our market vs traditional markets (on which there is no way to directly buy the underlying Silver)


We're selling future production at discount, a highly profitable margin for all parties involved.


\subsection{For owners, producers}

This process takes the assets out of circulation by digitizing it, the assets get locked on our chain, and most likely can no longer be easily acquired, so this increases scarcity in market


They get liquidity faster


As the digital equivalent trades, giving them new revenue stream from trading fees which increases lifetime value of each oz they extract


\subsection{For asset consumers}

Increased ability to time market and source adequate key resources


Ability to acquire resources at extreme discount


Ability to participate in market making, liquidity providing, digital ownership for increased revenue


\subsection{For society}

We're making key assets available cheaply and efficiently for everyone


We're able to enforce sustainable and environmentally friendly standards


We're able to enrich local communities around each mine with our digitization strategies


\section{Opportunity}

The Lux Silver project presents a unique opportunity in the intersection of precious metals and digital assets. By tokenizing real-world silver assets, Lux Silver combines the tangible value of silver with the flexibility and accessibility of blockchain technology. This is made possible through a strategic partnership with Viscount Mining, a reputable mining company with a portfolio of silver and gold properties in the Western United States. The Silver Cliff project, managed by Viscount Mining, is a high-grade, historic silver-gold mine located in Colorado, offering a unique and valuable asset for tokenization. The global silver market is robust and growing, driven by its use in various industries, including electronics, solar energy, and jewelry. Lux Silver, through its innovative approach and strong partnerships, provides exposure to this growing market in a flexible and accessible format. The forward Sales Agreement between Lux Partners Ltd. and Viscount Mining provides Lux with the opportunity to tokenize ten million ounces of silver, offering a unique and innovative investment opportunity. Furthermore, Lux Silver allows for the redemption of digital assets for physical silver, providing a level of security and tangibility often lacking in digital assets.


\section{About Lux}

Lux is a FinTech company domiciled in the Isle of Man and partnered with a regulated and licensed money transmitter business. Lux enables institutions to take advantage of blockchain technology use-cases in a tax-advantaged and regulated environment, with proper compliance, KYC, and AML procedures. It is therefore a secure and efficient solution to store, move, and grow assets at scale. Lux executives have an extensive track record of managing transactions and investments across a wide range of industries, which serves to provide assurance on the success of this transaction and support diligence efforts by leveraging in-house expertise and global networks.


\section{Viscount Mining Corp.}

Viscount Mining is a project generator with a portfolio of silver and gold properties in the Western United States, including Silver Cliff in Colorado and Cherry Creek in Nevada. The company's strategy is to identify and acquire properties that have substantial past production and historical resources and engage industry experts to mitigate execution risk and increase the probability of success.


The company's primary focus is the development of major mineral deposits in two areas in the USA, Colorado and Nevada. In Colorado, Viscount Mining holds a 100\% interest in the Silver Cliff property, a silver-gold project with a rich history and substantial silver resources. In Nevada, the company's Cherry Creek Property includes more than 400 unpatented and patented claims as well as mill rights, and the property includes more than 20 past producing mines.


Viscount Mining is committed to following the highest industry standards in all of its exploration and development activities, including adherence to environmental best practices and respect for the rights and cultures of local communities.


The company's partnership with Lux Partners Ltd. in the Silver Cliff project represents a significant milestone in the company's growth strategy. The forward Sales Agreement provides Lux with the opportunity to tokenize ten million ounces of silver, offering a unique and innovative investment opportunity in the silver market. This agreement, combined with the proven resources of the Silver Cliff project, positions Viscount Mining as a leader in the silver mining industry.


\section{Silver Cliff Project}

The Silver Cliff project is a high-grade, historic silver-gold mine located in the mining-friendly Colorado jurisdiction. Viscount Mining has 100\% ownership of the 2,029-acre property, which is easily accessible via county and state roads from the nearby major highway.


The Silver Cliff property lies in the historic Hardscrabble Silver District in the Wet Mountain Valley. Silver was discovered here in 1878, and by 1883, Silver Cliff was the third largest city in Colorado. The town boasted a population of over 5,000 people and was the county seat of Custer County. The Silver Cliff property has been the site of numerous historic mines. Silver Cliff ranks as one of the largest producers of silver in Colorado history.


The Silver Cliff project has a rich history of mining dating back to the late 19th century. The mine has already produced millions of ounces of silver and has records of more than 100 million ounces of silver at that location. The NI 43-101 report for the Silver Cliff project indicates 20 million ounces of silver. This report brings the historical records of the Silver Cliff mine up to the current NI 43-101 standard, providing a reliable estimate of the silver resources available.


The Silver Cliff property has excellent infrastructure. The property is less than an hour drive on state highways and county roads from Pueblo, Colorado, population 108,000. The property is on the electricity grid, and there is ample local accommodation for mine workers.


The Silver Cliff mine has a significant amount of silver resources, with the NI 43-101 report indicating 20 million ounces of silver. This report brings the historical records of the Silver Cliff mine up to the current NI 43-101 standard, providing a reliable estimate of the silver resources available.


\section{Tokenization Process}

Tokenization: Each ounce of silver is represented as an NFT (ERC-721 token), with a unique ID and metadata detailing the silver's specific characteristics. A Silver Token (LUXS), an ERC-20 token, is created as an algorithmic stablecoin backed by a basket of ERC-721s, enabling a dynamic, market dictated price based on the underlying silver reserves.


Pre-sale: Conduct a pre-sale of NFTs representing future silver. Each NFT is guaranteed to correspond to an ounce of silver. These NFTs are sold at a 20\% discount to the current spot price to incentivize buyers. This allows you to raise funds before the silver is even mined.


Liquidity Pool Launch: Create a LUXS/USDT liquidity pool on a decentralized exchange. This allows for the Silver Tokens to be traded against a stablecoin (USDT), ensuring liquidity and enabling price discovery for your tokens. In parallel to this we will enable users to dynamically set pricing for liquidity directly of various pools of NFTs as well, enabling direct liquidity for NFT holders that wish to sell, without going through physical redemption.


Trading: Once the liquidity pool is live, users can buy and sell Silver Tokens (ERC-20 and ERC-721) on the open market. Users who bought NFTs in the pre-sale can also trade them in secondary markets.


Mining and attaching Physical Silver to NFTs: When a pre-sold ounce of silver is mined and meets the specifications set in the NFT, "attach" the physical silver to its corresponding NFT. The silver's unique details are embedded in the NFT's metadata, and ownership of the NFT is transferred to the buyer. The physical silver remains in secure storage until the buyer chooses to take possession.


Redemption: Allow for the strict redemption of physical silver under set guidelines, such as minimum holding periods and/or thresholds. This can be enforced through the smart contracts underlying the NFTs. This way, you ensure that the majority of the tokens remain in circulation as long as possible, maintaining liquidity and market activity.


\section{Risk Factors}

Market volatility: The value of silver, like all precious metals, can fluctuate based on market conditions. This can affect the value of the Silver Tokens.


Regulatory changes: Changes in laws or regulations related to mining or blockchain technology could impact the project.


Mining risks: There are inherent risks associated with mining, including environmental risks, operational risks, and the risk that the actual quantity or quality of silver mined may be less than estimated.


\section{Market Outlook}

Silver, a precious metal with the highest reflectivity and the highest electrical and thermal conductivity of any metal, has been valued for its diverse applications since antiquity. It is used in various ways, such as in coins and bars as currency or as an investment, in solar panels, water filtration devices, in photographic and X-ray film, medical instruments, and in jewelry. Its ductility and high reflectivity make silver very attractive as jewelry, and its non-reactivity to oxygen and hydrogen prevent it from rusting like other metals.


Today, most silver is won in chemical processes after the extraction of lead ore and copper ore. Only about 20 percent of the silver that is produced today is extracted from silver ore reserves. In medieval times, most of the world's silver was produced in Europe, but in the nineteenth century the world's primary production of silver moved to North America. This still holds true, as Mexico is presently the global leader in silver mining.


In 2022, the global jewelry industry accounted for 234.1 million ounces of the global demand for silver, which was 18.8 percent of the total global silver demand that year. This indicates a robust demand for silver, not just for its industrial applications, but also for its aesthetic appeal in the form of jewelry.


The global silver market is robust and growing, driven by its use in various industries. The promise of future silver deliveries, combined with the novelty and potential investment value of NFTs and LUXS tokens, should drive strong sales, long before the silver is mined and processed. The forward Sales Agreement between Lux Partners Ltd. and Viscount Mining provides Lux with the opportunity to tokenize ten million ounces of silver, offering a unique and innovative investment opportunity.


Given these market conditions, the Lux Silver project presents a unique opportunity for investors to gain exposure to the silver market through a novel and flexible format. The tokenization of silver assets provides a level of liquidity and accessibility that is often lacking in traditional investment vehicles, making it an attractive option for a wide range of investors.


\section{Increased Market Demand from Solar Panels}

The photovoltaic (PV) industry's demand for silver is rapidly growing, with the International Energy Association (IEA) predicting solar investment to exceed oil production investment in 2023. The scarcity of polysilicon limited solar panel production in 2022, but by 2023, an overabundance of polysilicon has come online, causing panel input prices to fall and production to surge. Silver usage in solar panels was reduced by around 80\% over the past decade, but new solar technologies are likely to reverse this trend, requiring 30\% to 80\% more silver. The industry's production ramp-up continues to exceed prior expectations, with 350 GW in panel installations now seeming realistic for 2023, translating to at least 180 million ounces of silver needed. As solar panel production is increasing at astonishing rates and the switch to silver-rich TOPCon or HJT panels continues, we can expect an ongoing surge of physical silver demand from the photovoltaics industry. This demand cannot be satisfied by financial derivatives, nor can it be the kind of unallocated silver which is often sold without substantial backing for investors. Photovoltaic demand is poised to drain global physical reserves of silver. Should the 2024 production capacity reach 1,000 GW as predicted by IEA, the required photovoltaic silver demand would likely surpass 500 million troy ounces of silver by mid-decade. It is very unlikely that such demand could be satisfied without drastically higher silver prices or much high physical metal premiums. In conclusion, the photovoltaic industry's impending impact on silver pricing cannot be overlooked. As the silver shortages loom on the horizon, it is only a matter of time before mainstream media reports on these significant developments. Once the news spreads, we can anticipate substantial price surges in the silver market.


\end{document}
